\chapter{Normalización de Componentes y Recálculo del Punto Q}
\label{chap:normalizacion}

Los valores de resistencia obtenidos analíticamente en el capítulo previo corresponden a valores ideales continuos.
Para la implementación práctica en laboratorio se deben seleccionar los valores comerciales más próximos dentro de la
serie normalizada E12 (tolerancia nominal del $\pm 10\%$, comercialmente disponibles con tolerancia del $\pm 5\%$).
Posteriormente se recalcula el punto de reposo $Q$ para constatar que las desviaciones permanezcan dentro de los márgenes
admisibles fijados por la cátedra.

\section{Selección de Resistores Comerciales de la Serie E12}

Al normalizar la red de polarización, el objetivo primordial consiste en preservar con máxima fidelidad la tensión de
compuerta de reposo ($V_G \approx 3.294\V$) y el nivel de impedancia equivalente ($R_G \approx 100\kohm$). Se adoptan
los siguientes componentes de la serie comercial E12:
\begin{itemize}
  \item Resistor de surtidor: $R_S = 1.5\kohm$ (coincide con el valor calculado de $1.500\kohm$, error de $0.00\%$).
  \item Resistor de drenador: $R_D = 6.8\kohm$ (valor normalizado frente a $6.750\kohm$ calculados, desviación de $+0.74\%$).
  \item Resistor inferior de compuerta: $R_1 = 130\kohm$ (frente a $128.136\kohm$ calculados, desviación de $+1.45\%$).
  \item Resistor superior de compuerta: $R_2 = 470\kohm$ (frente a $455.420\kohm$ calculados, desviación de $+3.20\%$).
  \item Resistor de carga: $R_L = 10\kohm$ (componente estándar de serie E12).
\end{itemize}

La resistencia equivalente de compuerta con los valores comerciales adoptados resulta:
\begin{equation}
  R_{eq} = R_1 \parallel R_2 = \frac{130\kohm \cdot 470\kohm}{130\kohm + 470\kohm} = \frac{61100}{600}\kohm \approx 101.833\kohm
  \label{eq:Req_norm}
\end{equation}
con un desvío mínimo de apenas $+1.83\%$ respecto a los $100\kohm$ de diseño.

En el Cuadro~\ref{tab:resistores_normalizados} se resumen los valores calculados y normalizados adoptados.

\begin{table}[!ht]
  \centering
  \begin{tabular}{|l|c|c|c|c|c|}
    \hline
    \textbf{Componente} & $R_1$ & $R_2$ & $R_S$ & $R_D$ & $R_{eq} = R_1 \parallel R_2$ \\
    \hline
    Diseño Teórico $[\kohm]$ & $128.136$ & $455.420$ & $1.500$ & $6.750$ & $100.000$ \\
    Normalizado E12 $[\kohm]$ & $130.000$ & $470.000$ & $1.500$ & $6.800$ & $101.833$ \\
    Error Relativo $[\%]$ & $+1.45\%$ & $+3.20\%$ & $0.00\%$ & $+0.74\%$ & $+1.83\%$ \\
    \hline
  \end{tabular}
  \caption{Comparativa entre los resistores ideales de diseño y los componentes comerciales normalizados.}
  \label{tab:resistores_normalizados}
\end{table}

\section{Recálculo Analítico del Punto de Reposo}

Con los componentes comerciales adoptados, la nueva tensión continua del divisor de compuerta queda determinada por:
\begin{equation}
  V_G = V_{DD} \cdot \frac{R_1}{R_1 + R_2} = 15.0\V \cdot \frac{130\kohm}{130\kohm + 470\kohm} = 15.0\V \cdot \frac{130}{600} = 3.250\V
  \label{eq:VG_norm}
\end{equation}
presentando una variación de $-1.33\%$ respecto al valor teórico ($3.294\V$).

Para determinar la nueva corriente de drenador de reposo ($I_D$), se plantea la ecuación de la malla compuerta-surtidor:
\begin{equation}
  V_{GS} = V_G - I_D \cdot R_S
  \label{eq:malla_vgs}
\end{equation}
Sustituyendo en la ecuación cuadrática de saturación del MOSFET:
\begin{equation}
  I_D = \frac{K}{2} \cdot (V_G - I_D \cdot R_S - V_{th})^2
\end{equation}

Definiendo la variable auxiliar $x = V_G - V_{th} = 3.250\V - 1.60104\V = 1.64896\V \approx 1.649\V$, se reescribe:
\begin{equation}
  I_D = \frac{K}{2} \cdot (x - I_D \cdot R_S)^2 = \frac{K}{2} \cdot \left(x^2 - 2 \cdot x \cdot R_S \cdot I_D + R_S^2 \cdot I_D^2\right)
\end{equation}
Reordenando los términos algebraicamente en forma de ecuación cuadrática estándar $a \cdot I_D^2 + b \cdot I_D + c = 0$:
\begin{equation}
  \left(\frac{K \cdot R_S^2}{2}\right) \cdot I_D^2 - (1 + K \cdot x \cdot R_S) \cdot I_D + \left(\frac{K \cdot x^2}{2}\right) = 0
  \label{eq:cuadratica_ID}
\end{equation}

Evaluando con los parámetros numéricos normalizados:
\begin{equation*}
  K = 53.905\mA/\V^2 = 0.053905\A/\V^2, \quad R_S = 1500\ohm, \quad x = 1.64896\V
\end{equation*}
\begin{equation}
  a = \frac{0.053905 \cdot 1500^2}{2} = 60643.125\ \Omega^2\A/\V^2
\end{equation}
\begin{equation}
  b = -(1 + 0.053905 \cdot 1.64896 \cdot 1500) = -134.330
\end{equation}
\begin{equation}
  c = \frac{0.053905 \cdot 1.64896^2}{2} = 0.073285\ \A
\end{equation}

Resolviendo la ecuación de segundo grado mediante la fórmula resolvente:
\begin{equation}
  I_D = \frac{134.330 - \sqrt{134.330^2 - 4 \cdot 60643.125 \cdot 0.073285}}{2 \cdot 60643.125} \approx 0.9727\mA
  \label{eq:ID_norm_val}
\end{equation}
donde se descarta la raíz complementaria por carecer de significado físico ($V_{GS} < V_{th}$).

A partir de este valor de corriente se deducen secuencialmente los restantes potenciales nodales de continua:
\begin{equation}
  V_S = I_D \cdot R_S = 0.9727\mA \cdot 1.5\kohm = 1.459\V
  \label{eq:VS_norm_val}
\end{equation}
\begin{equation}
  V_D = V_{DD} - I_D \cdot R_D = 15.0\V - (0.9727\mA \cdot 6.8\kohm) = 15.0\V - 6.614\V = 8.386\V
  \label{eq:VD_norm_val}
\end{equation}
\begin{equation}
  V_{GS} = V_G - V_S = 3.250\V - 1.459\V = 1.791\V
  \label{eq:VGS_norm_val}
\end{equation}
\begin{equation}
  V_{DS} = V_D - V_S = 8.386\V - 1.459\V = 6.927\V
  \label{eq:VDS_norm_val}
\end{equation}

Se verifica nuevamente la condición de saturación:
\begin{equation}
  V_{DS} = 6.927\V > V_{GS} - V_{th} = 1.791\V - 1.601\V = 0.190\V
\end{equation}
cumpliéndose holgadamente con la premisa de diseño.

En el Cuadro~\ref{tab:recalculo_punto_q} se contrastan los valores del punto de trabajo estático teórico y normalizado.

\begin{table}[!ht]
  \centering
  \resizebox{0.88\textwidth}{!}{
  \begin{tabular}{|l|c|c|c|}
    \hline
    \textbf{Parámetro Eléctrico} & \textbf{Teórico (Diseño)} & \textbf{Con R Normalizadas} & \textbf{Error Relativo} \\
    \hline
    Tensión de compuerta $V_G\ [\V]$ & $3.294$ & $3.250$ & $-1.33\%$ \\
    Exceso de compuerta $V_G - V_{th}\ [\V]$ & $1.693$ & $1.649$ & $-2.58\%$ \\
    Corriente de drenador $I_D\ [\mA]$ & $1.000$ & $0.9727$ & $-2.73\%$ \\
    Tensión de surtidor $V_S\ [\V]$ & $1.500$ & $1.459$ & $-2.73\%$ \\
    Tensión de drenador $V_D\ [\V]$ & $8.250$ & $8.386$ & $+1.65\%$ \\
    Tensión compuerta-surtidor $V_{GS}\ [\V]$ & $1.794$ & $1.791$ & $-0.15\%$ \\
    Tensión drenador-surtidor $V_{DS}\ [\V]$ & $6.750$ & $6.927$ & $+2.62\%$ \\
    \hline
  \end{tabular}
  }
  \caption{Comparación entre el punto de reposo analítico teórico y el obtenido con componentes normalizados.}
  \label{tab:recalculo_punto_q}
\end{table}

Se observa que la mayor discrepancia relativa es de apenas $-2.73\%$ en la corriente de drenador y tensión de surtidor,
lo cual comprueba la excelente robustez del esquema de polarización adoptado frente a la dispersión de componentes
comerciales de la serie E12.
